\chapter{Subversion Repository Hook Reference}\label{svn.ref.reposhooks}

Subversion repositories provide a number of event hooks which are
essentially opportunities for administrators to extend
Subversion\textquotesingle s functionality at key moments of key
operations. Repository hooks are implemented as programs executed by
Subversion itself at those key moments---before and after a commit,
before and after a user locks a file, and so on.

For each hook it provides, Subversion will attempt to execute the
program of that hook\textquotesingle s name which is found in the
\texttt{hooks/} subdirectory of the repository\textquotesingle s on-disk
directory structure. For example, on a Unix system, the start-commit
hook script would be installed at
\texttt{REPOS\_PATH/hooks/start-commit}, where it could be a binary
executable program, a shell script, a Python program, etc. On a Windows
system, the program would be installed in the same location, but would
be named \texttt{START-COMMIT.EXE} or \texttt{START-COMMIT.BAT} instead
of simply \texttt{start-commit}.

This reference guide describes the various hooks which Subversion offers
to administrators, detailing when the hook is invoked, its input
parameters, and how its behavior affects the Subversion workflow.

\section{start-commit}\label{svn.ref.reposhooks.start-commit}

\emph{Notification of the beginning of a commit.}

\textbf{Synopsis}

\texttt{start-commit\ REPOS-PATH\ USER\ CAPABILITIES\ TXN-NAME}

\subsection{Description}

The start-commit hook is run immediately after the commit transaction is
created and its initial properties set. It is typically used as an early
termination mechanism, avoiding what could be a lengthy commit process
which would eventually fail at a later phase anyway due to a
user\textquotesingle s lack of commit privileges or some other commit
metadata validation failure.

If the start-commit hook program returns a nonzero exit value, the
commit process is stopped, the commit transaction is destroyed, and
anything printed to \texttt{stderr} is marshalled back to the client.

The start-commit hook is not a suitable place to evaluate the substance
of a particular commit, as it is invoked before any file or directory
change information has been transmitted. Use the pre-commit hook script
(which is described in
\hyperref[svn.ref.reposhooks.pre-commit]{pre-commit} elsewhere in this
reference) for that purpose.

Prior to Subversion 1.8, the Subversion invoked the start-commit hook
\emph{before} creating the commit transaction. Failure of the script
resulted in that transaction not being created at all. This was changed
in Subversion 1.8, though, to allow implementations of the start-commit
hook access to the transaction\textquotesingle s properties, which can
include (among other things) the revision log associated with the
commit.

\subsection{Input Parameter(s)}

The command-line arguments passed to the hook program, in order, are:

\begin{enumerate}
\def\labelenumi{\arabic{enumi}.}
\item
  Repository path
\item
  Authenticated username attempting the commit
\item
  Colon-separated list of capabilities that a client passes to the
  server, including \texttt{depth}, \texttt{mergeinfo}, and
  \texttt{log-revprops} (new in Subversion 1.5)
\item
  Commit transaction name (new in Subversion 1.8)
\end{enumerate}

\subsection{Common uses}

Access control (e.g., temporarily lock out commits for some reason).

A means to allow access only from clients that have certain
capabilities.

\section{pre-commit}\label{svn.ref.reposhooks.pre-commit}

\emph{Notification just prior to commit completion.}

\textbf{Synopsis}

\texttt{pre-commit\ REPOS-PATH\ TXN-NAME}

\subsection{Description}

The pre-commit hook is run just before a commit transaction is promoted
to a new revision. Typically, this hook is used to protect against
commits that are disallowed due to content or location (e.g., your site
might require that all commits to a certain branch include a ticket
number from the bug tracker, or that the incoming log message is
nonempty).

If the pre-commit hook program returns a nonzero exit value, the commit
is aborted, the commit transaction is removed, and anything printed to
\texttt{stderr} is marshalled back to the client.

\subsection{Input parameter(s)}

The command-line arguments passed to the hook program, in order, are:

\begin{enumerate}
\def\labelenumi{\arabic{enumi}.}
\item
  Repository path
\item
  Commit transaction name
\end{enumerate}

Additionally, Subversion passes any lock tokens provided by the
committing client to the hook script via standard input. When present,
these are formatted as a single line containing the string
\texttt{LOCK-TOKENS:}, followed by additional lines---one per lock
token---which contain the lock token information. Each lock token
information line consists of the URI-escaped repository filesystem path
associated with the lock, followed by the pipe (\texttt{\textbar{}})
separator character, and finally the lock token string.

\subsection{Common uses}

Change validation and control

\section{post-commit}\label{svn.ref.reposhooks.post-commit}

\emph{Notification of a successful commit.}

\textbf{Synopsis}

\texttt{post-commit\ REPOS-PATH\ REVISION\ TXN-NAME}

\subsection{Description}

The post-commit hook is run after the transaction is committed and a new
revision is created. Most people use this hook to send out descriptive
emails about the commit or to notify some other tool (such as an issue
tracker) that a commit has happened. Some configurations also use this
hook to trigger backup processes.

If the post-commit hook returns a nonzero exit status, the commit
\emph{will not} be aborted since it has already completed. However,
anything that the hook printed to \texttt{stderr} will be marshalled
back to the client, making it easier to diagnose hook failures.

\subsection{Input parameter(s)}

The command-line arguments passed to the hook program, in order, are:

\begin{enumerate}
\def\labelenumi{\arabic{enumi}.}
\item
  Repository path
\item
  Revision number created by the commit
\item
  Name of the transaction that has become the revision triggering the
  post-commit hook.
\end{enumerate}

\subsection{Common uses}

Commit notification; tool integration

\section{pre-revprop-change}\label{svn.ref.reposhooks.pre-revprop-change}

\emph{Notification of a revision property change attempt.}

\textbf{Synopsis}

\texttt{pre-revprop-change\ REPOS-PATH\ REVISION\ USER\ PROPNAME\ ACTION}

\subsection{Description}

The pre-revprop-change hook is run immediately prior to the modification
of a revision property when performed outside the scope of a normal
commit. Unlike the other hooks, the default state of this one is to deny
the proposed action. The hook must actually exist and return a zero exit
value before a revision property modification can happen.

If the pre-revprop-change hook doesn\textquotesingle t exist,
isn\textquotesingle t executable, or returns a nonzero exit value, no
change to the property will be made, and anything printed to
\texttt{stderr} is marshalled back to the client.

\subsection{Input parameter(s)}

The command-line arguments passed to the hook program, in order, are:

\begin{enumerate}
\def\labelenumi{\arabic{enumi}.}
\item
  Repository path
\item
  Revision whose property is about to be modified
\item
  Authenticated username attempting the property change
\item
  Name of the property changed
\item
  Change description: \texttt{A} (added), \texttt{D} (deleted), or
  \texttt{M} (modified)
\end{enumerate}

Additionally, Subversion passes the intended new value of the property
to the hook program via standard input.

\subsection{Common uses}

Access control; change validation and control

\section{post-revprop-change}\label{svn.ref.reposhooks.post-revprop-change}

\emph{Notification of a successful revision property change.}

\textbf{Synopsis}

\texttt{post-revprop-change\ REPOS-PATH\ REVISION\ USER\ PROPNAME\ ACTION}

\subsection{Description}

The post-revprop-change hook is run immediately after the modification
of a revision property when performed outside the scope of a normal
commit. As you can derive from the description of its counterpart, the
pre-revprop-change hook, this hook will not run at all unless the
pre-revprop-change hook is implemented. It is typically used to send
email notification of the property change.

If the post-revprop-change hook returns a nonzero exit status, the
change \emph{will not} be aborted since it has already completed.
However, anything that the hook printed to \texttt{stderr} will be
marshalled back to the client, making it easier to diagnose hook
failures.

\subsection{Input parameter(s)}

The command-line arguments passed to the hook program, in order, are:

\begin{enumerate}
\def\labelenumi{\arabic{enumi}.}
\item
  Repository path
\item
  Revision whose property was modified
\item
  Authenticated username of the person making the change
\item
  Name of the property changed
\item
  Change description: \texttt{A} (added), \texttt{D} (deleted), or
  \texttt{M} (modified)
\end{enumerate}

Additionally, Subversion passes to the hook program, via standard input,
the previous value of the property.

\subsection{Common uses}

Property change notification

\section{pre-lock}\label{svn.ref.reposhooks.pre-lock}

\emph{Notification of a path lock attempt.}

\textbf{Synopsis}

\texttt{pre-lock\ REPOS-PATH\ PATH\ USER\ COMMENT\ STEAL}

\subsection{Description}

The pre-lock hook runs whenever someone attempts to lock a path. It can
be used to prevent locks altogether or to create a more complex policy
specifying exactly which users are allowed to lock particular paths. If
the hook notices a preexisting lock, it can also decide whether a user
is allowed to ``steal'' the existing lock.

If the pre-lock hook program returns a nonzero exit value, the lock
action is aborted and anything printed to \texttt{stderr} is marshalled
back to the client.

The hook program may optionally dictate the lock token which will be
assigned to the lock by printing the desired lock token to standard
output. Because of this, implementations of this hook should carefully
avoid unexpected output sent to standard output.

If the pre-lock script takes advantage of lock token dictation feature,
the responsibility of generating a \emph{unique} lock token falls to the
script itself. Failure to generate unique lock tokens may result in
undefined---and very likely, undesired---behavior.

\subsection{Input parameter(s)}

The command-line arguments passed to the hook program, in order, are:

\begin{enumerate}
\def\labelenumi{\arabic{enumi}.}
\item
  Repository path
\item
  Versioned path that is to be locked
\item
  Authenticated username of the person attempting the lock
\item
  Comment provided when the lock was created
\item
  \texttt{1} if the user is attempting to steal an existing lock;
  \texttt{0} otherwise
\end{enumerate}

\subsection{Common uses}

Access control

\section{post-lock}\label{svn.ref.reposhooks.post-lock}

\emph{Notification of a successful path lock.}

\textbf{Synopsis}

\texttt{post-lock\ REPOS-PATH\ USER}

\subsection{Description}

The post-lock hook runs after one or more paths have been locked. It is
typically used to send email notification of the lock event.

If the post-lock hook returns a nonzero exit status, the lock \emph{will
not} be aborted since it has already completed. However, anything that
the hook printed to \texttt{stderr} will be marshalled back to the
client, making it easier to diagnose hook failures.

\subsection{Input parameter(s)}

The command-line arguments passed to the hook program, in order, are:

\begin{enumerate}
\def\labelenumi{\arabic{enumi}.}
\item
  Repository path
\item
  Authenticated username of the person who locked the paths
\end{enumerate}

Additionally, the list of paths locked is passed to the hook program via
standard input, one path per line.

\subsection{Common uses}

Lock notification

\section{pre-unlock}\label{svn.ref.reposhooks.pre-unlock}

\emph{Notification of a path unlock attempt.}

\textbf{Synopsis}

\texttt{pre-unlock\ REPOS-PATH\ PATH\ USER\ TOKEN\ BREAK-UNLOCK}

\subsection{Description}

The pre-unlock hook runs whenever someone attempts to remove a lock on a
file. It can be used to create policies that specify which users are
allowed to unlock particular paths. It\textquotesingle s particularly
important for determining policies about lock breakage. If user A locks
a file, is user B allowed to break the lock? What if the lock is more
than a week old? These sorts of things can be decided and enforced by
the hook.

If the pre-unlock hook program returns a nonzero exit value, the unlock
action is aborted and anything printed to \texttt{stderr} is marshalled
back to the client.

\subsection{Input parameter(s)}

The command-line arguments passed to the hook program, in order, are:

\begin{enumerate}
\def\labelenumi{\arabic{enumi}.}
\item
  Repository path
\item
  Versioned path which is to be unlocked
\item
  Authenticated username of the person attempting the unlock
\item
  Lock token associated with the lock which is to be removed
\item
  \texttt{1} if the user is attempting to break the lock; \texttt{0}
  otherwise
\end{enumerate}

\subsection{Common uses}

Access control

\section{post-unlock}\label{svn.ref.reposhooks.post-unlock}

\emph{Notification of a successful path unlock.}

\textbf{Synopsis}

\texttt{post-unlock\ REPOS-PATH\ USER}

\subsection{Description}

The post-unlock hook runs after one or more paths have been unlocked. It
is typically used to send email notification of the unlock event.

If the post-unlock hook returns a nonzero exit status, the unlock
\emph{will not} be aborted since it has already completed. However,
anything that the hook printed to \texttt{stderr} will be marshalled
back to the client, making it easier to diagnose hook failures.

\subsection{Input parameter(s)}

The command-line arguments passed to the hook program, in order, are:

\begin{enumerate}
\def\labelenumi{\arabic{enumi}.}
\item
  Repository path
\item
  Authenticated username of the person who unlocked the paths
\end{enumerate}

Additionally, the list of paths unlocked is passed to the hook program
via standard input, one path per line.

\subsection{Common uses}

Unlock notification

